\conclusions

\chapter*{CONCLUSIONES FINALES}
\addcontentsline{toc}{chapter}{CONCLUSIONES FINALES}

Tras finalizar el proceso de investigación, diseño e implementación de la plataforma Assets Studio, se han obtenido los siguientes resultados que dan cumplimiento a los objetivos planteados:

\begin{enumerate}
    \item La sistematización de los referentes teóricos permitió identificar que la gestión de activos digitales (DAM) moderna debe trascender el simple almacenamiento, integrando principios de experiencia de usuario (UX) y arquitecturas escalables. Se confirmó que, para el contexto cubano, es imperativo adaptar estas tecnologías a condiciones de baja conectividad mediante optimizaciones en el rendimiento y el uso de infraestructuras \textit{cloud-native}.

    \item El diagnóstico del estado actual evidenció una problemática crítica en el sector creativo nacional: una gestión de recursos fragmentada y desorganizada que genera una pérdida de hasta el 30\% del tiempo laboral en tareas logísticas. Esta ineficiencia, sumada a la dispersión de los activos en múltiples dispositivos y la falta de herramientas de búsqueda inteligente, justificó plenamente la necesidad de desarrollar una solución centralizada.

    \item El diseño y la implementación de Assets Studio demostraron la viabilidad de utilizar un \textit{stack} tecnológico moderno basado en Next.js, PostgreSQL y servicios en la nube como Cloudinary y Firebase. La integración de capacidades de Inteligencia Artificial para el etiquetado automático y la búsqueda semántica resultó ser el núcleo transformador de la solución, permitiendo pasar de una organización manual basada en carpetas a una gestión contextual y eficiente.

    \item La validación de la propuesta confirmó la alta efectividad de la solución, obteniendo una calificación de 85/100 en la escala de usabilidad (SUS). Las pruebas empíricas demostraron una reducción drástica en los tiempos de localización de activos (de minutos a pocos segundos) y una mejora significativa en la colaboración de los equipos. Asimismo, el estudio de factibilidad ratificó que la plataforma es económicamente sostenible y técnicamente robusta para su despliegue inmediato.
\end{enumerate}